\section{Comparación de métodos y compensaciones de diseño}

\subsection{Comparación de tres aproximaciones a la función de verosimilitud}

\begin{table}[H]
\centering
\begin{tabular}{p{0.2\textwidth}p{0.25\textwidth}p{0.2\textwidth}p{0.2\textwidth}}
\toprule
\textbf{Método} & \textbf{Aproximación de $p(x_0|x_t)$} & \textbf{Requisito para $\mathcal{A}$} & \textbf{Precisión de la aproximación} \\
\midrule
Aproximación más rudimentaria & $p(y|x_t) = p(y|x_0)$ & Cualquiera & Peor \\
DPS & $\delta(x_0 - \hat{x}_0)$ & Diferenciable es suficiente & Inferior \\
Pseudo-inversa guiada & $\mathcal{N}(\hat{x}_0, r_t^2 I)$ & Debe ser lineal & Relativamente mejor \\
\bottomrule
\end{tabular}
\caption{Comparación de tres métodos de aproximación a la función de verosimilitud}
\end{table}

\begin{importantbox}{Compensación entre precisión de la aproximación y alcance de aplicación}
Estos tres métodos ilustran una compensación de diseño central:
\begin{itemize}
    \item Las aproximaciones más precisas (como la pseudo-inversa guiada) requieren supuestos más fuertes ($\mathcal{A}$ debe ser lineal), pero son teóricamente más confiables.
    \item Las aproximaciones más rudimentarias (como DPS) tienen requisitos menores para $\mathcal{A}$ (solo necesita ser diferenciable), tienen un alcance de aplicación más amplio, pero la error de aproximación es mayor.
\end{itemize}
No existe ningún método que sea óptimo en todas las circunstancias; la elección depende del caso de uso específico y de las propiedades del mapeo hacia adelante.
\end{importantbox}

\subsection{Limitaciones comunes de todos los métodos}

\begin{warningbox}{Limitaciones inherentes de la guía por aproximación}
Independientemente de la aproximación utilizada, estos métodos de guía durante la inferencia comparten las siguientes limitaciones:
\begin{enumerate}
    \item \textbf{No garantizan convergencia a la distribución posterior correcta}: Todas las aproximaciones introducen errores; los resultados generados pueden no satisfacer completamente las condiciones de restricción.
    \item \textbf{Requieren ajuste de hiperparámetros}: La intensidad de la guía (step size / escala de guía) necesita ser ajustada manualmente, y diferentes tareas y conjuntos de datos pueden requerir valores distintos.
    \item \textbf{Costo computacional}: DPS requiere retropropagación del estimador de Tweedie, lo que incrementa el costo computacional por paso.
    \item \textbf{Las aproximaciones en pasos de tiempo con alto ruido son peores}: En las etapas tempranas (con alto ruido), la calidad de la estimación de $\hat{x}_0$ es muy pobre, lo que hace que la señal de guía sea poco confiable.
\end{enumerate}
\end{warningbox}

\subsection{Resumen de la sección}

Diferentes métodos de aproximación a la función de verosimilitud forman un espectro entre precisión y alcance de aplicación. En aplicaciones prácticas, es necesario seleccionar el método adecuado basándose en las propiedades del mapeo hacia adelante, el presupuesto computacional y los requisitos de calidad de resultados.


